% !TEX root = ../tesi.tex
\chapter{Stato dell'Arte}
\label{chap:stato_arte}

L'applicazione della computer vision in agricoltura ha conosciuto una rapida 
evoluzione negli ultimi anni, passando da semplici algoritmi basati sul colore a 
modelli avanzati di deep learning capaci di interpretare scene complesse. 
Questo capitolo esamina lo stato dell'arte delle tecnologie di visione per la raccolta 
robotizzata, partendo dall'evoluzione dei modelli di rilevamento e segmentazione 
fino ad analizzare le tecniche proposte in letteratura per affrontare il problema 
delle occlusioni nei grappoli densi.


\section{Computer Vision per l'Agricoltura di Precisione}
\label{sec:cv_agricoltura}

I primi tentativi di applicare la visione artificiale alla raccolta automatizzata 
risalgono a diversi decenni fa e si basavano su tecniche classiche di elaborazione 
delle immagini \cite{gongal2015sensors}. Questi approcci tradizionali utilizzavano 
principalmente la segmentazione per soglia cromatica (ad esempio negli spazi colore 
RGB o HSV), filtri morfologici per pulire il rumore di fondo e semplici algoritmi 
geometrici per l'analisi dei contorni.

Tuttavia, questi metodi mostravano limiti evidenti nelle condizioni reali di una serra. 
La segmentazione basata sul colore funzionava solo in laboratorio sotto luce artificiale 
stabile, mentre all'interno della serra la luce solare cambia continuamente durante la 
giornata, mandando in crisi le soglie fisse. Il problema principale riguardava la forte 
somiglianza cromatica tra i pomodori acerbi e la vegetazione circostante: le foglie e i 
frutti verdi condividono tonalità quasi identiche, rendendo impossibile distinguerli 
affidandosi al solo colore. Inoltre, le tecniche tradizionali non erano in grado di 
gestire le frequenti sovrapposizioni tra frutti, finendo spesso per considerare interi 
grappoli come un'unica sagoma indistinta.

La svolta è arrivata con l'introduzione del \textit{Deep Learning} e in particolare 
delle reti neurali convoluzionali (CNN). A differenza degli approcci classici, in cui le 
regole per riconoscere un frutto dovevano essere scritte a mano, le reti neurali sono 
capaci di apprendere autonomamente le caratteristiche visive a partire dalle immagini 
fornite durante l'addestramento. Il modello impara così a combinare informazioni di colore, 
bordi, sfumature e contesto della scena. Questo cambio di paradigma ha permesso di 
riconoscere i pomodori anche in presenza di forti variazioni di luce, ombre fitte e 
grappoli molto densi, gettando le basi percettive per i moderni sistemi di raccolta robotizzata.

\section{Rilevamento e Segmentazione delle Istanze}
\label{sec:detection_segmentation}

All'interno delle tecniche di deep learning per la visione artificiale, i modelli 
per la localizzazione degli oggetti si dividono storicamente in due grandi categorie: 
gli approcci a due stadi (come la famiglia R-CNN \cite{he2017mask}), che prima generano 
regioni candidate e poi le classificano, e i modelli a singolo stadio (\textit{single-stage}), 
tra cui spicca l'architettura YOLO (\textit{You Only Look Once}) \cite{redmon2016you}.

I modelli YOLO hanno introdotto una vera e propria rivoluzione nel settore: invece di 
spezzare l'elaborazione in fasi separate, una singola rete neurale elabora l'immagine 
in un unico passaggio, producendo direttamente le coordinate degli oggetti e le relative 
classi. Questa caratteristica garantisce una velocità di esecuzione molto elevata, 
rendendo l'algoritmo adatto a funzionare in tempo reale a bordo di sistemi robotici.

Nel campo della raccolta robotizzata, è fondamentale distinguere tra due diversi 
compiti di visione che i modelli YOLO moderni possono svolgere:
\begin{itemize}
    \item \textbf{Object Detection}: il modello localizza ciascun oggetto presente 
    nella scena racchiudendolo all'interno di una \textit{bounding box} rettangolare, 
    a cui associa una classe e un punteggio di confidenza. Questo approccio è molto rapido, 
    ma per sua natura approssimato: il rettangolo indica la presenza del pomodoro ma 
    include anche porzioni di sfondo, rami e foglie, senza fornire alcuna informazione 
    sulla sagoma reale del frutto.
    \item \textbf{Instance Segmentation}: rappresenta un'estensione più avanzata, in cui 
    il modello non si limita a disegnare un rettangolo, ma assegna un'etichetta a ogni 
    singolo pixel, producendo una maschera binaria che ritaglia la forma esatta di ciascun 
    pomodoro \cite{liu2020yolotomato}. In un'applicazione di raccolta robotica, disporre della 
    maschera a livello di pixel è indispensabile: consente di distinguere pomodori a stretto 
    contatto, calcolarne con precisione il centroide e delimitare i margini fisici necessari 
    per posizionare correttamente l'organo di presa.
\end{itemize}

All'interno del progetto di ricerca del nostro laboratorio, la prima fase del lavoro 
è stata dedicata proprio allo studio, all'addestramento e all'ottimizzazione di questi 
modelli percettivi su un dataset acquisito in serra \cite{collega_yolo}. Da tale indagine 
comparativa è emerso che il modello \textbf{YOLO11-x} (nella sua configurazione per 
l'Instance Segmentation) garantisce il miglior compromesso tra accuratezza metrica e 
qualità delle maschere prodotte, riuscendo a distinguere con efficacia i pomodori anche 
in presenza di fogliame fitto.

Tuttavia, disporre delle maschere di segmentazione fornite da YOLO costituisce solo il 
primo passo. La rete neurale si limita a dire "dove si trovano i pomodori", ma non fornisce 
alcuna informazione su come essi interagiscano nello spazio tridimensionale, chi copra chi 
e quale debba essere l'ordine di raccolta. È esattamente su questo divario che si innesta 
il presente lavoro di tesi.


\section{Tecniche di Gestione delle Occlusioni in Letteratura}
\label{sec:occlusioni_letteratura}

La gestione delle occlusioni reciproche all'interno di grappoli densi rappresenta 
uno degli aspetti più complessi nella robotica agricola \cite{bac2014harvesting}. 
Quando i frutti crescono a stretto contatto, determinare quale pomodoro si trovi 
davanti agli altri e stabilire un ordine corretto di raccolta è indispensabile per 
consentire all'organo di presa di operare in sicurezza senza rischiare di danneggiare 
i frutti circostanti \cite{silwal2017design}.

All'interno dell'architettura complessiva sviluppata nel nostro laboratorio, l'obiettivo 
applicativo a lungo termine prevede l'impiego di una telecamera RGB-D montata direttamente 
sul braccio meccanico (\textit{eye-in-hand}), così da guidare il manipolatore verso il target 
combinando l'informazione visiva con la stima della distanza \cite{collega_yolo}. 
Tuttavia, i dati raccolti sul campo in serra e messi a disposizione per la fase di addestramento 
dei modelli percettivi consistevano esclusivamente in immagini a colori bidimensionali (RGB). 
Inoltre, la letteratura scientifica evidenzia come l'elaborazione diretta di nuvole di punti 
3D in ambiente agricolo presenti sensibili criticità: i sensori di profondità risentono delle 
variazioni di illuminazione solare tipiche delle serre, la vegetazione sottile introduce 
rumore geometrico e gli algoritmi di ricostruzione volumetrica richiedono un carico computazionale 
elevato, difficilmente sostenibile durante i cicli di controllo rapido del braccio 
\cite{gongal2015sensors, bac2014harvesting}. Si è quindi reso necessario concentrare questo 
lavoro di tesi sullo sviluppo di un modulo decisionale che operi interamente sul piano 
bidimensionale delle immagini RGB, estraendo la gerarchia spaziale delle occlusioni prima 
di inviare il punto di mira all'attuatore.

Sul piano bidimensionale, il primissimo approccio che abbiamo sviluppato insieme al collega 
nella fase iniziale del lavoro — subito dopo aver completato l'addestramento del modello 
YOLO — è consistito nel valutare le occlusioni sfruttando direttamente i riquadri di 
delimitazione rettangolari (\textit{bounding box}) forniti nativamente dalla rete, 
in linea con diversi lavori preliminari presenti in letteratura \cite{chen2011practical}. 
Tuttavia, l'applicazione pratica di questa prima soluzione ha mostrato fin da subito un grave 
limite geometrico: i pomodori hanno una conformazione curvilinea, per cui i quattro angoli della 
scatola rettangolare racchiudono una percentuale consistente di spazio vuoto. Quando due 
frutti si trovano vicini o disposti lungo linee diagonali, le rispettive bounding box finiscono 
per sovrapporsi vistosamente anche se i corpi dei pomodori non sono affatto a contatto. 
Questo fenomeno genera falsi blocchi geometrici, inducendo il sistema a considerare occlusi 
pomodori che nella realtà sarebbero perfettamente liberi e accessibili.

In letteratura, alcuni studi hanno proposto di ridurre queste imprecisioni ricorrendo 
all'involucro convesso (\textit{Convex Hull}) per regolarizzare i contorni dei frutti 
\cite{changhui2017reconstruction} o all'approssimazione mediante ellissi orientate 
\cite{wang2024multifruit}. Tuttavia, la maggior parte di queste indagini si è concentrata 
principalmente sulla ricostruzione morfologica del singolo frutto isolato, senza sviluppare 
un algoritmo decisionale completo per stabilire la sequenza ottima di prelievo.

Un ulteriore limite comune a gran parte dei metodi esistenti risiede nella natura statica 
della pianificazione. Nella quasi totalità dei casi, l'ordine di raccolta viene calcolato 
in un unico passaggio iniziale, ordinando i pomodori in base alla sola visibilità o alla distanza. 
Questo approccio ignora la dinamica fisica dell'operazione di raccolta: il prelievo del primo 
pomodoro frontale modifica immediatamente la scena, liberando lo spazio di manovra e rendendo 
accessibili i frutti retrostanti che in precedenza risultavano bloccati.

Il presente lavoro di tesi nasce esattamente dalla necessità di superare il nostro approccio 
iniziale a \textit{bounding box} e le limitazioni dei metodi statici. L'obiettivo è proporre 
una pipeline 2D leggera e deterministica che, operando sulle maschere di segmentazione di 
YOLO11-x, utilizzi il Convex Hull per ripulire le imperfezioni, modelli i frutti tramite ellissi 
orientate e introduca un algoritmo dinamico (\textit{Pick \& Update}) capace di ricalcolare le 
priorità a ogni prelievo simulato. Nel capitolo successivo verrà descritta nel dettaglio 
la formulazione matematica e l'architettura completa di questa metodologia.


